\subsection*{Datasets}
The planned primary benchmark is URPC2018, which provides four common underwater biological
categories. To evaluate generalization, we also consider RUOD as a larger and more diverse
benchmark with more classes and more complex scenes.

\subsection*{Implementation Details}
The initial baseline is planned around a YOLO-style detector. We will report results on a
Linux environment with a complete CUDA stack and an RTX 3090 GPU. Evaluation metrics will
include mAP@0.5, mAP@0.5:0.95, precision, recall, parameter count, GFLOPs, and FPS.

\subsection*{Comparison Experiments}
The final paper will compare QDCR-Net against standard lightweight detectors and recent
underwater-oriented baselines. The comparison will emphasize low-visibility performance,
small-object detection accuracy, and efficiency trade-offs.

\subsection*{Ablation Studies}
The ablation plan includes:
\begin{enumerate}[leftmargin=*]
  \item baseline detector;
  \item enhanced-input-only detector;
  \item dual branch without cross residual interaction;
  \item dual branch with cross residual interaction;
  \item dual branch with cross residual interaction and quality-aware fusion;
  \item full model with the small-object neck.
\end{enumerate}

Additional mechanism analysis will compare concat/add fusion against cross-residual design,
single-direction versus bidirectional interaction, fixed versus dynamic fusion, and different
insertion stages for the interaction block.

\subsection*{Visualization and Robustness Analysis}
We will include qualitative detection maps, failure cases, and grouped robustness studies
under blur, low contrast, and color-cast conditions. Small-object subsets and low-visibility
subsets will be evaluated separately to verify the intended gains of the method.
